\documentclass[11pt,a4paper]{article}
% preamble.tex — estilo común de todos los apuntes Froth.
% Cada apunte hace % preamble.tex — estilo común de todos los apuntes Froth.
% Cada apunte hace % preamble.tex — estilo común de todos los apuntes Froth.
% Cada apunte hace \input{preamble} justo después de \documentclass.
\usepackage[margin=2.4cm]{geometry}
\usepackage{xcolor}
\definecolor{litblue}{RGB}{31,79,121}
\definecolor{litgreen}{RGB}{27,94,32}
\definecolor{litorange}{RGB}{230,81,0}
\usepackage{parskip}
\usepackage{titlesec}
\titleformat{\section}{\Large\bfseries\color{litblue}}{\thesection}{0.6em}{}
\titleformat{\subsection}{\large\bfseries\color{litblue}}{\thesubsection}{0.6em}{}
\usepackage{tcolorbox}
\usepackage{fancyhdr}
\pagestyle{fancy}
\fancyhf{}
\fancyhead[L]{\small\color{litblue}Froth — Apuntes de ML}
\fancyhead[R]{\small\thepage}
\renewcommand{\headrulewidth}{0.4pt}
\usepackage[colorlinks=true,linkcolor=litblue,urlcolor=litblue]{hyperref}

% Cabecera de cada apunte: \cabecera{numero}{titulo}
\newcommand{\cabecera}[2]{%
  \begin{tcolorbox}[colback=litblue,coltext=white,colframe=litblue,arc=2mm]
    {\Large\bfseries Apunte #1 — #2}\\[3pt]
    {\small Proyecto Froth · aprender Machine Learning construyendo}
  \end{tcolorbox}\vspace{0.8em}}

% Cajas temáticas
\newtcolorbox{concepto}[1]{colback=blue!4,colframe=litblue,title={Concepto: #1},arc=1.5mm}
\newtcolorbox{practica}{colback=green!5,colframe=litgreen,title={Pruébalo tú (teclado en mano)},arc=1.5mm}
\newtcolorbox{ojo}{colback=orange!8,colframe=litorange,title={Ojo / error típico},arc=1.5mm}
\newtcolorbox{resumen}{colback=gray!8,colframe=gray!60,title={En una frase},arc=1.5mm}

\newcommand{\codigo}[1]{\texttt{#1}}
 justo después de \documentclass.
\usepackage[margin=2.4cm]{geometry}
\usepackage{xcolor}
\definecolor{litblue}{RGB}{31,79,121}
\definecolor{litgreen}{RGB}{27,94,32}
\definecolor{litorange}{RGB}{230,81,0}
\usepackage{parskip}
\usepackage{titlesec}
\titleformat{\section}{\Large\bfseries\color{litblue}}{\thesection}{0.6em}{}
\titleformat{\subsection}{\large\bfseries\color{litblue}}{\thesubsection}{0.6em}{}
\usepackage{tcolorbox}
\usepackage{fancyhdr}
\pagestyle{fancy}
\fancyhf{}
\fancyhead[L]{\small\color{litblue}Froth — Apuntes de ML}
\fancyhead[R]{\small\thepage}
\renewcommand{\headrulewidth}{0.4pt}
\usepackage[colorlinks=true,linkcolor=litblue,urlcolor=litblue]{hyperref}

% Cabecera de cada apunte: \cabecera{numero}{titulo}
\newcommand{\cabecera}[2]{%
  \begin{tcolorbox}[colback=litblue,coltext=white,colframe=litblue,arc=2mm]
    {\Large\bfseries Apunte #1 — #2}\\[3pt]
    {\small Proyecto Froth · aprender Machine Learning construyendo}
  \end{tcolorbox}\vspace{0.8em}}

% Cajas temáticas
\newtcolorbox{concepto}[1]{colback=blue!4,colframe=litblue,title={Concepto: #1},arc=1.5mm}
\newtcolorbox{practica}{colback=green!5,colframe=litgreen,title={Pruébalo tú (teclado en mano)},arc=1.5mm}
\newtcolorbox{ojo}{colback=orange!8,colframe=litorange,title={Ojo / error típico},arc=1.5mm}
\newtcolorbox{resumen}{colback=gray!8,colframe=gray!60,title={En una frase},arc=1.5mm}

\newcommand{\codigo}[1]{\texttt{#1}}
 justo después de \documentclass.
\usepackage[margin=2.4cm]{geometry}
\usepackage{xcolor}
\definecolor{litblue}{RGB}{31,79,121}
\definecolor{litgreen}{RGB}{27,94,32}
\definecolor{litorange}{RGB}{230,81,0}
\usepackage{parskip}
\usepackage{titlesec}
\titleformat{\section}{\Large\bfseries\color{litblue}}{\thesection}{0.6em}{}
\titleformat{\subsection}{\large\bfseries\color{litblue}}{\thesubsection}{0.6em}{}
\usepackage{tcolorbox}
\usepackage{fancyhdr}
\pagestyle{fancy}
\fancyhf{}
\fancyhead[L]{\small\color{litblue}Froth — Apuntes de ML}
\fancyhead[R]{\small\thepage}
\renewcommand{\headrulewidth}{0.4pt}
\usepackage[colorlinks=true,linkcolor=litblue,urlcolor=litblue]{hyperref}

% Cabecera de cada apunte: \cabecera{numero}{titulo}
\newcommand{\cabecera}[2]{%
  \begin{tcolorbox}[colback=litblue,coltext=white,colframe=litblue,arc=2mm]
    {\Large\bfseries Apunte #1 — #2}\\[3pt]
    {\small Proyecto Froth · aprender Machine Learning construyendo}
  \end{tcolorbox}\vspace{0.8em}}

% Cajas temáticas
\newtcolorbox{concepto}[1]{colback=blue!4,colframe=litblue,title={Concepto: #1},arc=1.5mm}
\newtcolorbox{practica}{colback=green!5,colframe=litgreen,title={Pruébalo tú (teclado en mano)},arc=1.5mm}
\newtcolorbox{ojo}{colback=orange!8,colframe=litorange,title={Ojo / error típico},arc=1.5mm}
\newtcolorbox{resumen}{colback=gray!8,colframe=gray!60,title={En una frase},arc=1.5mm}

\newcommand{\codigo}[1]{\texttt{#1}}

\begin{document}
\cabecera{09}{El buscador semántico (juntamos todo)}

\section{Qué construimos}

Una herramienta que recibe una \textbf{frase en lenguaje natural} y devuelve los papers más
parecidos \emph{por significado} --- no por palabras compartidas. Es la función
\codigo{most\_similar()} de \codigo{froth/embed.py}, y une todo lo aprendido en la Fase 2.

\begin{concepto}{búsqueda semántica vs.\ búsqueda por palabras}
La búsqueda clásica (Ctrl+F, Google básico) empareja \textbf{palabras}: si buscas ``grain
size'' no encuentra un texto que diga ``coarse and fine particles'', aunque signifiquen lo
mismo. La búsqueda \textbf{semántica} empareja \textbf{significado}: convierte tu frase en un
vector y busca los vectores cercanos, así que sí los relaciona.
\end{concepto}

\section{Cómo funciona por dentro (3 pasos)}

\begin{enumerate}
  \item \textbf{Vectorizar la frase.} Se pasa tu frase por el \emph{mismo} modelo (SPECTER)
        con \codigo{model.encode([query])}. Clave: la frase y los papers deben vivir en el
        mismo mapa, o comparar no tendría sentido.
  \item \textbf{Medir contra todos.} Con \codigo{cosine\_similarity} calculamos el coseno de
        tu frase contra los 126 vectores de golpe.
  \item \textbf{Quedarnos con los mejores.} Ordenamos de mayor a menor coseno y devolvemos los
        \codigo{k} primeros (por defecto 5), como una tabla \codigo{score, title, year}.
\end{enumerate}

\begin{ojo}
El mismo modelo para las dos cosas. Si vectorizaras los papers con SPECTER y la frase con
otro modelo distinto, los dos ``mapas'' no coincidirían y los resultados serían basura.
Consistencia del embedding = regla de oro.
\end{ojo}

\section{La prueba (resultados reales)}

Búsqueda: \emph{``how does grain size affect the recovery of minerals by froth flotation''}
\begin{center}
\begin{tabular}{ll}
\textbf{Score} & \textbf{Paper} \\
\hline
0.873 & Role of Bubble Size in Flotation of Coarse and Fine Particles \\
0.856 & Effect of Clay Slime on the Froth Stability and Flotation Performance \\
0.855 & Spodumene Flotation Mechanism \\
\end{tabular}
\end{center}
Escribimos ``grain size'' y encontró papers de ``coarse and fine particles'' y ``ultrafine'':
\textbf{entendió el significado, no las palabras}.

Búsqueda: \emph{``extracting beryllium and tantalum as valuable by-products''}
\begin{center}
\begin{tabular}{ll}
\textbf{Score} & \textbf{Paper} \\
\hline
0.812 & Niobium and tantalum \\
0.740 & Concentration of Tantalum and Niobium \\
\end{tabular}
\end{center}
Trajo justo los papers de los by-products Ta/Nb/Be de tu tesis.

\begin{concepto}{score entre 0 y 1}
El \codigo{score} es la similitud coseno (Apunte 08). Aquí ronda 0.7--0.9: alto, tiene
sentido. Ojo: los scores de ``frase libre vs.\ paper'' suelen ser algo más bajos que
``paper vs.\ paper'' (0.94 antes), porque una frase corta es menos rica que un abstract
entero. Lo que importa es el \emph{orden}, no el número absoluto.
\end{concepto}

\section{Qué significa esto para el proyecto}

Ya tienes una herramienta útil de verdad: interrogar tu corpus en lenguaje natural. Y más
importante: la maquinaria (vectorizar + coseno + ordenar) es \textbf{la misma} que en la
Fase 3 formará las burbujas de subtemas. Dominar esto es dominar el núcleo de Froth.

\begin{resumen}
El buscador semántico convierte tu frase en vector, la compara por coseno contra los 126
papers y devuelve los más parecidos por significado. Probado: ``grain size'' encontró
``fine/coarse particles''. El motor de Froth ya late.
\end{resumen}

\end{document}
